\documentclass[10pt,a4paper]{amsart}
\usepackage[margin=19mm]{geometry}
\usepackage{amsmath,amssymb,mathtools}
\usepackage[scr=rsfs,cal=euler]{mathalfa}
\usepackage{xfrac}
\usepackage[unicode,hidelinks,bookmarks=false]{hyperref}
\newcommand{\ie}{i.e.\ }
\newcommand{\cf}{cf.\ }
\newcommand{\resp}{respectively\ }
\newcommand{\Subsection}{\subsection}
\providecommand{\nobreakdash}{\nobreak\hbox{-}}
\setlength{\emergencystretch}{2em}
\multlinegap0pt
\usepackage{mathtools}
\usepackage{xcolor}
\usepackage{amssymb}
\usepackage{bm}
\usepackage{tikz}
\newtheorem{lem}{Lemma}
\usepackage{cleveref}
\crefname{lem}{Lemma}{Lemmas}
\newcommand{\Sh}{\mathrm{Sh}}
\newcommand{\Z}{\mathbb{Z}}
\newcommand{\ZhatP}{\widehat{\Z}_\mc{P}}
\newcommand{\colim}{\operatorname{\mathrm{colim}}}
\newcommand{\et}{\mathrm{\acute{e}t}}
\newcommand{\mc}{\mathcal}
\newcommand{\proet}{\mathrm{pro\acute{e}t}}
\title{Pro-étale motives and solid rigidity}
\author{Rapha\"el Ruimy, Swann Tubach, Sebastian Wolf}
\date{Source: arXiv:2601.07358v1 (CC BY 4.0)}
\begin{document}
\maketitle

\noindent\emph{Source and licence.} Pro-étale motives and solid rigidity. Version arXiv:2601.07358v1 (CC BY 4.0), distributed by arXiv under the Creative Commons Attribution 4.0 International licence, http://creativecommons.org/licenses/by/4.0/, as recorded in the arXiv licence metadata for this exact version and shown on the abstract page https://arxiv.org/abs/2601.07358v1. Full licence notice: \texttt{licenses/arxiv-2601.07358-CC-BY-4.0.txt}.\par\smallskip

\noindent\textit{Editorial scope.} Counted unit: the article's lem lem:solidification-of-general-pro-etale-map (source0.tex lines 1068--1077). The article's complete original proof of it is delivered with it (source0.tex lines 1078--1104, one \texttt{proof} environment of 27 lines). No mathematics is written, repaired or re-derived: the statement and the proof are the article's own lines, in the article's own notation, and no retained line is substituted or re-flowed.  Retained uncounted as the source-local context the proof needs: lines 873--877.  Nothing was omitted from any retained span, so the package carries no omission record.  No retained line contains a citation, so the wrapper carries no bibliography and no bibliography entry.  No retained line contains a figure environment, an image-inclusion command or drawing code, so no image is delivered and the package declares no asset dependency.  Editorial formatting only, in addition: an AMS \texttt{amsart} preamble that restates the article's own declarations and macros used by the retained lines, namely the environments \texttt{lem} on separate counters, together with the standard packages those lines need.  The retained environments print in this document as Lemma 1 in order of appearance, whereas the article prints the same items with the numbering of its own full document.  The complete native source of the article as distributed in the arXiv e-print for this exact version is bundled unchanged as \texttt{sources/192494/source0.tex}.
\begin{lem}
\label{limIsRlim}
	Let $(F_i)$ be a cofiltered inverse system of
$\mc{P}$-torsion constructible étale sheaves. Then for all $q>0$, the $R^q\lim F_i$ vanishes when taken in the category $\Sh(X_\proet,\ZhatP)$. %of pro-étale sheaves on $X$.
\end{lem}

\begin{lem}
	\label{lem:solidification-of-general-pro-etale-map}
	Let $V = \lim_{\alpha \in A} V_\alpha$ be an inverse limit of qcqs \'etale $X$-schemes with affine transition maps and $j \colon V \to X$, $j_\alpha \colon V_\alpha \to X$ the structure maps.
	Then the canonical map
	\[
		j_\sharp^\blacksquare\ZhatP\coloneqq (j_\sharp \ZhatP)^{\blacksquare} \to \lim_\alpha (j_\alpha)_\sharp \ZhatP
	\]
	is an equivalence in $\Sh(X,\ZhatP)^{\blacksquare}$.
	In particular the full subcategory $\Sh(X,\ZhatP)^{\blacksquare} \subseteq \Sh(X_\proet,\ZhatP)$ is equivalent to the reflective localization at all maps $j_\sharp\ZhatP \to j_\sharp^\blacksquare\ZhatP\simeq\lim_\alpha (j_\alpha)_\sharp \ZhatP$. 
\end{lem}

\begin{proof}
	The second part of the statement is immediate from the first, so we only show the first part.
	For this recall that by \Cref{limIsRlim} the functor \begin{equation*}\mathrm{Pro}(\Sh_\mc{P}^\mathrm{cons}(X_\et)) \to \Sh(X_\proet,\ZhatP)\end{equation*} is exact so in particular it preserves finite colimits.
	Without loss of generality we may assume that the indexing poset admits a final object $\alpha_0$.
	Since $V_{\alpha_0}$ is qcqs, we can write $V_{\alpha_0} = \colim_i V_{\alpha_0}^i$ as a finite colimit where the $V_{\alpha_0}^i$ are affine \'etale $X$-schemes.
	It follows that we have
	\[
		V_{\alpha} \cong \colim_i V_{\alpha}^i
	\]
	where $V_{\alpha}^i =  V_{\alpha_0}^i \times_{V_{\alpha_0}}  V_{\alpha}$ for any $\alpha \in A$.
	Since taking free $\ZhatP$-sheaves preserves colimits, we therefore also get
	\[
		(j_\alpha)_\sharp \ZhatP \cong \colim_i (j_\alpha^i)_\sharp \ZhatP
	\] with $j_\alpha^i\colon V_\alpha^i\to X$ the natural map.
	Since finite colimits commute with inverse limits in $\mathrm{Pro}(\Sh_\mc{P}^\mathrm{cons}(X_\et))$ it follows that
	\[
		\lim_\alpha(\colim_i (j_\alpha^i)_\sharp \ZhatP) \cong \colim_i (\lim_\alpha (j_\alpha^i)_\sharp \ZhatP)
	\]
	in $\Sh(X_\proet,\ZhatP)$.
	Let us write $j^i \colon V^i = \lim_\alpha V_{\alpha}^i \to X$ for the projection. 
	Then one shows similarly as above that
	\[
		j_\sharp \ZhatP \cong \colim_i j^i_\sharp \ZhatP.
	\]
	Thus to show that $(j_\sharp \ZhatP)^{\blacksquare} \to \lim_\alpha (j_\alpha)_\sharp \ZhatP$ is an equivalence it suffices to see that $(j^i_\sharp \ZhatP)^{\blacksquare} \to \lim_\alpha (j^i_\alpha)_\sharp \ZhatP$ is an equivalence for all $i$.
	This is clear since now all the schemes involved are affine.
\end{proof}

\end{document}
